

%------------------------------------------------------------------
 % Title: sublinearity of LSH based on p-stable mapping

%Introduction
% Preliminaries: how the IM98 algo works and relation to LSH
% Effectively a new LSH for l2
% the algorithm : p-stable distribution
% the analysis in terms of a general pdf
% plots for l1 and l2 and improvement over 1/c
% analysis for p < 1.
% conclusion

\newcommand{\remove}[1]{}

\documentclass[12pt]{article}

% All style files are available from 
%   http://wwww.uiuc.edu/~sariel/research/latex/
\usepackage[breaklinks]{hyperref}  % comment out if you do not have it
\usepackage{graphicx}
\usepackage[active]{srcltx} % comment out if you do not have it
\usepackage{sariel,wide}

%\addtolength{\topmargin}{-1.0cm}
%\addtolength{\textheight}{2.4cm}
%\addtolength{\oddsidemargin}{-0.75cm}
%\addtolength{\evensidemargin}{-0.75cm}
%\addtolength{\textwidth}{1.5cm}

%\newcommand{\llabel}[1]{~\fbox{\scriptsize{\sf #1}}\label{#1}}
\newcommand{\llabel}[1]{\label{#1}}
\newcommand{\tw}{{\rm tw}}
\newcommand{\ltw}{{\rm ltw}}
%\newtheorem{fact}{Fact}

\newcommand{\bm}[1]{\mbox{\boldmath$#1$}}

%\newcommand{\eps}{\epsilon}
\newcommand{\MS}{{\cal M}}

%\newcommand{\remove}[1]{}
%\newcommand{\remove}[1]{#1}
\newcommand{\eop}{\hfill$\Box$}
\renewcommand{\H}{\end{document}}
%\def\baselinestretch{.97}
%\mainmatter
\pagestyle{plain}
\begin{document}

%\newtheorem{theorem}{Theorem}
%\newtheorem{definition}{Definition}
%\newtheorem{lemma}{Lemma}

%\remove{}
%\setcounter{page}{0}

\title{Locality-Sensitive Hashing Scheme Based on p-Stable Distributions}

\author{Mayur Datar\thanks{Department of Computer Science,; Stanford
      University; {\tt datar@cs.stanford.edu}}%
   \and%
   Nicole Immorlica\thanks{Laboratory for Computer Science, MIT;
      {\tt nickle@theory.lcs.mit.edu}} %
   \and%
   Piotr Indyk\thanks{Laboratory for Computer Science, MIT; {\tt
         indyk@theory.lcs.mit.edu} }
   \and Vahab S. Mirrokni\thanks{Laboratory for Computer Science, MIT;
   {\tt mirrokni@theory.lcs.mit.edu}}}

\date{}

%\begin{titlepage}
\maketitle

\begin{abstract}
We present a novel Locality-Sensitive Hashing scheme for the Approximate 
Nearest Neighbor Problem under $l_p$ norm, based on $p$-stable distributions.

Our scheme improves the running time of the earlier algorithm for the case of the $l_2$ norm.
It also yields the first known provably efficient approximate NN algorithm for 
the case $p <1$.  We also show that the algorithm finds the {\em exact} near neigbhor in $O(\log n)$ time for data satisfying certain ``bounded growth'' condition.

Unlike earlier schemes, our LSH scheme works directly on points 
in the Euclidean space without embeddings. Consequently, the resulting 
query time bound is free of large factors and is simple and easy to 
implement. Our experiments (on synthetic data sets) show that the our data structure is up to 40 
times faster than $kd$-tree.
\end{abstract}

\iffalse
\vspace{2mm}
\noindent {\bf Topic area:} core database technology

\noindent {\bf Special Purpose DB Techn.:} Multidimensional Databases, Advanced Search, Query, and Approximation

\noindent {\bf Paper ID:} 620
\fi
%\end{titlepage}

\section{Introduction}
A similarity search problem involves a collection of objects (documents,
images, etc.) that are characterized by a collection of relevant features
and represented as points in a high-dimensional attribute space; given
queries in the form of points in this space, we are required to find
the nearest (most similar) object to the query.
A particularly interesting and well-studied instance is 
$d$-dimensional Euclidean space.
% under some $l_p$ norm.
This problem
is of major importance to a variety of applications;
% usually involving
%similarity searching.
some examples are:
data compression,
databases and data mining,
information retrieval,
image and video databases,
machine learning, 
pattern recognition,  statistics and data analysis.
Typically, the features of the objects of interest (documents, images, etc)
are represented as points in $\Re^d$ and a distance metric is used to measure
similarity of objects. The basic problem then is to perform indexing or
similarity searching for query objects.
The number of features (i.e., the
dimensionality) ranges anywhere from tens to thousands. 
\iffalse
For example, in
multimedia applications such as IBM's QBIC (Query by Image Content),
the number of features could be several hundreds~\cite{fqbic,qbic}.
In information
retrieval for text documents, vector-space representations involve
several thousands of dimensions, and it is often considered to be
a dramatic improvement that the dimension-reduction techniques 
%such as
%LSI (latent semantic indexing)~\cite{BDS,BSMS,DDLFH},
%principal components analysis~\cite{hotelling}
%or the Karhunen-Lo\'eve transform~\cite{karhunen,loeve},
can reduce the dimensionality to a mere few hundreds!

%For $d=100$, we need $n$ larger than $2^{100}$ to make any of
%the known algorithms appear better than brute-force search. 

%Of late, there has been an increasing interest in avoiding the curse
%of dimensionality by resorting to {\em approximate} nearest neighbor
%searching~\cite{AM,Cl2,AMNSW,Cl3,Kl}. 
\fi

The low-dimensional case (say, for the dimensionality $d$ equal to 2 or 3) is well-solved,
so the main issue is that of dealing with a large number of dimensions,
the so-called ``curse of dimensionality''. Despite decades
of intensive effort, the current solutions are not entirely satisfactory;
in fact, for large enough $d$, in theory or in practice, they often provide little 
improvement over a linear algorithm which compares a query to each point from the database.
In particular, it was shown in~\cite{Web} (both empirically and theoretically) that {\em all} current indexing techniques (based on space partitioning) degrade to linear search for sufficiently high dimensions.

In recent years, several researchers   proposed to avoid the running time bottleneck by using {\em approximation} (e.g.,~\cite{AMNSW,Kl,IM,KOR,HP}, see 
also~\cite{NIPS}).
This is due to the fact that, in many cases, approximate nearest neighbor is almost as good as the exact one; in particular, if the distance measure accurately captures the notion of user quality, then small differences in the distance should not matter.
%Since the selection of features and the choice of a distance metric in the
%applications is rather heuristic and merely an attempt to make mathematically
%precise what is after all an essentially aesthetic notion of similarity,
%it is unnecessary to insist on the absolute nearest neighbor.
%In such cases, determining an {\em approximate} nearest neighbor,
%for a reasonable approximation value, should suffice for most practical
%purposes.
In fact, in situations when the quality of the approximate nearest neighbor is much worse than the quality of  the actual nearest neighbor,
then the nearest neighbor problem is {\em unstable}, and it is not clear if solving it is at all meaningful~\cite{BGRS,HAK}.

\iffalse
In addition, most of the approximate algorithms can be easily modified to report {\em all} approximate neighbors, including the exact one.
Since it is quite common (e.g., see Figure 3 in~\cite{GIM}) that the number of, say, $2$-approximate nearest neighbors (see definition in the next section) is small, an efficient 
approximate algorithm can report all such neighbors very quickly and then return the closest (exact) one.

%
%In most applications, it is much more important to keep the query time
%down to sub-linear in the data size, and a natural way to achieve this goal is to allow the algorithm to output only
%approximate nearest neighbors.
\fi



%What in many applications cannot be compromised is the query time.
%Imagine for example a search engine which enables content-based image
%retrieval on the World-Wide Web. If the system was to index a significant
%fraction of the web, the number of images to index would be at least of
%the order tens (if not hundreds) of million.
%In order to be able to answer a few million queries per day (as it is currently the case 
%for major {\em text-retrieval}  engines), the query time should be of the order of tens
%%of microseconds.
%Clearly, no indexing method exhibiting linear (or close to linear) dependence
%on the data size could manage such a huge data set.



In~\cite{IM,GIM}, the authors introduced an approximate high-dimensional similarity search scheme with provably sublinear dependence on the data size.
Instead of using tree-like space partitioning, it relied on a new method called
{\em locality-sensitive hashing (LSH)}.
The key idea is to hash the points using several hash functions so as to
ensure that, for each function, the probability of collision is much higher for
objects which are close to each other than for those which are far apart.
Then, one can determine near neighbors by hashing the query point and
retrieving elements stored in buckets containing that point.
In~\cite{IM,GIM} the authors provided such locality-sensitive hash functions for the case when the points live in binary Hamming space $\{0,1\}^d$.
They showed experimentally that the data structure achieves large speedup over several tree-based data structures when the data is stored on disk.
%
\iffalse
 that are simple and easy to implement;
they can also be naturally extended to the {\em dynamic} setting, i.e., when insertion and deletion operations also need to be supported.
%Another advantage of our scheme is that (as it is based on hashing) it can be naturally extended to the
%{\em dynamic} setting, i.e., when insertion and deletion operations also
%need to be supported.
\fi
%
In addition,  since the LSH is a hashing-based scheme, 
it can be naturally extended to the {\em dynamic} setting, i.e., when insertion and deletion operations also need to be supported.
This avoids the complexity of dealing with tree structures when the data is dynamic.


The LSH algorithm has been since used in numerous applied settings, 
e.g., see~\cite{GIM,eight,HGI,Shiv,Buh,BT,Ya,Buh2,OMST,GSM}.
However, it suffers from a fundamental drawback:
it is fast and simple only when the input points live in the Hamming space (indeed, almost all of the above applications  involved binary data).
As mentioned in~\cite{IM,GIM}, it is possible to extend the algorithm to the $l_2$ norm,
by embedding $l_2$ space into $l_1$ space, and then $l_1$ space into Hamming space.
However, it increases the query time and/or error by a large factor and complicates the algorithm.
\iffalse
Although the computation time is not crucial when the data is stored on disk 
(as the I/O cost dominates the running time), 
it is of crucial importance when the data is stored in main memory.
The latter scenario occurs more and more often, as memory is cheap and easily available.
\fi

In this paper we present a novel version of the LSH algorithm.
As with the previous schemes, it works for the $(R,c)$-Near Neighbor (NN) problem, where the goal is to report a point within distance $cR$ from a query $q$, if there is a point in the data set $P$ within distance $R$ from $q$.
Unlike the earlier algorithm, our algorithm works directly on points in Euclidean space without embeddings. 
As a consequence, it has the following advantages over the previous algorithm:
\begin{itemize}
\item For the $l_2$ norm, its query time is $O(d n^{\rho(c)}
\log n)$, where $\rho(c)<1/c$  for $c \in (1,10]$ (the inequality is strict, see
Figure~\ref{fig:l2_plot}). Thus, for large range of values of $c$, the query time exponent is better than the one in~\cite{IM,GIM}.
\item It is simple and quite easy to implement.
\item It works for any $l_p$ norm, as long as $p \in (0,2]$. Specifically, we show that for any $p \in (0,2]$ and $\gamma>0$ there exists an algorithm for $(R,c)$-NN under $l_d^p$
which uses  $O(d n +  n^{1+\rho})$ space, with query time 
$O(n^{\rho} \log_{1/\gamma} n)$, where where $\rho\le (1+\gamma)\cdot\max\left(\frac{1}{c^p}, \frac{1}{c}\right)$.
To our knowledge, this is the only known {\em provable} algorithm for the high-dimensional nearest neighbor problem for the case $p<1$.
Similarity search under such {\em fractional} norms have recently attracted interest~\cite{AHK,CIKM}.
\end{itemize}


Our algorithm also inherits two very convenient properties of LSH schemes.
The first one is that it works well on data that is extremely high-dimensional but sparse.
Specifically, the running time bound remains unchanged if $d$ denotes the maximum number of  non-zero elements in vectors.
To our knowledge, this property is not shared by other known spatial data structures.
Thanks to this property, we were able to use our new LSH scheme (specifically, the $l_1$ norm version) for fast color-based image similarity search~\cite{IT}.
In that context, each image was represented by a point in roughly $100^3$-dimensional space, but only about 100 dimensions were non-zero per point.
The use of our LSH scheme enabled achieving order(s) of magnitude speed-up over the linear scan.

The second property is that our algorithm provably reports the {\em exact} near neighbor very quickly, if the data satisfies certain {\em bounded growth} property.
Specifically, for a query point $q$, and $c \ge 1$,  let $N(q,c)$ be the number of $c$-approximate nearest neighbors of $q$ in $P$. 
If $N(q,c)$ grows ``sub-exponentially'' as a function of $c$, then the
LSH algorithm reports $p$, the nearest neighbor, with constant probability within time $O(d \log n)$, assuming it is given a constant factor approximation to the distance from $q$ to its nearest neighbor.
In particular, we show that if $N(q,c)=O(c^b)$, then the running time is $O(\log n +2^{O(b)})$. 
Efficient nearest neighbor algorithms for data sets with polynomial growth properties in general metrics have been recently a focus of several papers~\cite{Cl3,KRu,KLee}.
LSH solves an easier problem (near neighbor under $l_2$ norm), while working under weaker assumptions about the growth function.
It is also somewhat faster, due to the fact that the 
$\log n$ factor in the query time of the earlier schemes is {\em multiplied} by a function of $b$, while in our case this factor is additive. 

\iffalse
The main contribution of the paper is a LSH scheme for the $(R,c)$-Point
Location in Equal Balls (NN) problem\footnote{See Preliminaries for
definition} in $\Re^d$ under $l_p$ norm for $p \in (0,2]$. To the best
of our knowledge no LSH scheme was known earlier for $\Re^d$ under the
$l_p$ norm. It may be of independent interest as a sketching technique
for massive data sets, similar to the similarity estimation techniques
introduced by Charikar\cite{Char}.
\fi

We complement our theoretical analysis with experimental evaluation of the algorithm on data with wide range of parameters.
In particular, we compare our algorithm to an approximate version of 
the $kd$-tree algorithm~\cite{AMexp}.
We performed the experiments on synthetic data sets containing ``planted'' near neighbor (see section~\ref{Sec:expt} for more details); similar model was earlier used in~\cite{Yia2}.
Our experiments indicate that the new LSH scheme achieves query time of up to 40 times better than the query time of the $kd$-tree algorithm.


\subsection{Notations and problem definitions}
We use $l^d_p$ to denote the space $\Re^d$ under the $l_p$ norm. For any point
$\bm{v} \in \Re^d$, we denote by $||\vec{\bm{v}}||_p$ the $l_p$ norm of the vector
$\vec{\bm{v}}$. Let $\MS=(X,d)$ be any metric space, and $v \in X$. The {\em ball} of radius $r$ centered at $v$ is defined as
$B(v,r) = \{ q\in X \mid d(v,q) \leq r\}$.
%\subsection{The Near Neighbor formulation}

\iffalse
In this paper we focus on the $(R,\epsilon)$-Point Location in Equal
Balls (NN) problem.

\begin{definition}[$(R,\epsilon)$-NN]
\vspace*{-0.1in}
Given $n$ radius-$R$ balls centered at
$C=\{ c_1, \ldots, c_n \}$ in $\MS=(X,d)$, devise a data structure which for
any query point $q \in X$ does the following:
\vspace*{-0.1in}
\begin{itemize}
\item if there exists $c_i \in C$ with $q \in B(c_i,R)$ then return {\sc yes} and a point $c'_i$ such that $q \in B(c'_i,(1+\epsilon)R)$,
\item if $q \notin B(c_i,(1+\epsilon)R)$ for all $c_i \in C$ then return {\sc no},
\item if for the point $c_i$ closest to $q$ we have
$R \le d(q,c_i) \le ((1+\epsilon)R$) then return either {\sc yes} or {\sc no}.
\end{itemize}
\end{definition}
\fi

Let $c=1+\epsilon$.
In this paper we focus on the $(R,c)$-NN problem.
Observe that $(R,c)$-NN is simply a decision version of the Approximate Nearest Neighbor problem.
Although in many applications solving the decision version is good enough, one can also reduce the approximate NN problem to approximate NN via binary-search-like approach.
In particular, it is known~\cite{IM,HP} that
the $c$-approximate NN problem reduces  to
%\footnote{The original reduction from \cite{IM}
%was different and less efficient.  This was recently improved in \cite{HP}}
$O(\log (n/\epsilon))$ instances of
$(R,c)$-NN.
Then, the complexity of $c$-approximate NN is the
same (within log factor) as that of the $(R,c)$-NN problem.


\iffalse

\subsection{Formal statement of theoretical results}
\label{ss:formal}

%The main theoretical contribution of this paper is an analysis of an algorithm for solving the $(R,c)$-NN problem in $l_p^d$. 
The formal guarantees for our algorithm is given by the following theorem.

\begin{theorem}
For any $p \in (0,2]$ and $\gamma>0$ there exists an algorithm for $(R,c)$-NN under $l_d^p$
which uses  $O(d n +  n^{1+\rho})$ space, with query time 
$O(n^{\rho} \log_{1/\gamma} n)$
%evaluations of the hash function for each query,
where $\rho\le (1+\gamma)\cdot\max\left(\frac{1}{c^p}, \frac{1}{c}\right)$.
\end{theorem}


The proof of this theorem is omitted due to lack of space. 
We mention that it is somewhat challenging, due to the fact that there is no explicit formula that describes the density of $p$-stable distributions. Instead we use known properties of $p$-stable distributions (such as tail estimates, existence of first moments, etc.).

Instead we provide (for $p=1,2$) a computational analysis of the exponent 
$\rho=\rho(c)$ for $c \in [1,20]$.
See section~\ref{s:comp} for details.

\fi

\iffalse
It
admits the first known $\epsilon$-approximate NN algorithm for $p <1$. In
fact, for the special case of $p = 2$ (Euclidean norm) we show that the
exponent $\rho$ is strictly smaller than $1/c$ over a large range of
$c$. Thus, for this case it improves the previous best space and time
bounds (\cite{IM}) that follow from embedding $l_2^d$ into a Hamming space
(via embedding into $l_1^{O(d)}$).
\fi

\iffalse
\subsection{Previous work}


The nearest neighbor and its approximate versions are among the most extensively studied problems 
in the fields of databases, computational geometry and algorithms, resulting in the
discovery of
many  efficient algorithms.
In particular, for the case when the underlying metric space  is a low-dimensional
Euclidean space $l_2^d$, 
it is known how to construct data structures for the exact~\cite{Cl1,Me} or
the approximate~\cite{AMNSW,Kl,IM,KOR,HP} NN problem.
Unfortunately, those data structures have query time or space requirements
that are exponential in $d$. This is a severe bottleneck, since most
massive data sets (text documents, market-basket data, images etc.) that
need to be analyzed contain data points from a high
dimensional space. 
More recently, several data structures using space (quasi)-polynomial in $n$ and $d$, 
and query time 
sublinear in $n$, have been developed for approximate NN under $l_1$ and
$l_2$ ~\cite{KOR,IM} and $l_{\infty}$~\cite{I98} norms.
However, most of those algorithms are only of theoretical interest, 
since either the required space is too large, or the query time is significantly 
sublinear in $n$ only for very large approximation factors.
The only exception is the Locality-Sensitive Hashing (LSH) algorithm of~\cite{IM}, 
that solves the $c$-approximate NN problem (more specifically, its {\em decision} version).


There is considerable literature on the nearest neighbor problem in the database literature.
Samet~\cite{samet} surveys a variety of data structures
for nearest neighbors in geometric spaces, including variants of
$kd$-trees, $R$-trees, and structures based on space-filling curves.
The more recent results are surveyed in~\cite{SRF}; see also an excellent survey of~\cite{BK2}. 
Those algorithms often exhibit exponential dependence on dimension in the worst case and
usually in typical cases as well (e.g., see~\cite{AMN}).
In general, even the average-case
analysis of heuristics for points distributed over regions in $\Re^d$
gives an exponential query time~\cite{bentley,FBF,samet}.
\fi

\section{Locality-sensitive hashing}
\label{subsec:lsh}

An important technique from \cite{IM}, to solve the
$(R,c)$-NN problem is locality sensitive hashing or LSH.
For a domain $S$ of the points set with distance measure $D$, an LSH family
is defined as:
\begin{defn}
\label{defn:lsh}
A family ${\cal H} = \{h: S \to U\}$ is called
$(r_1,r_2,p_1,p_2)${\em-sensitive} for $D$  if for any $v, q \in S$
\vspace*{-0.1in}
\begin{itemize}
\item if $v \in B(q,r_1)$ then $\Pr_{\cal H}[h(q)=h(v)] \ge p_1$,
\item if $v \notin B(q,r_2)$ then $\Pr_{\cal H}[h(q)=h(v)] \le p_2$.
\end{itemize}
\end{defn}
In order for a locality-sensitive hash (LSH) family to be useful, it has to satisfy
inequalities $p_1>p_2$ and $r_1<r_2$.
% when $D$ is a dissimilarity measure,
%or $p_1>p_2$ and $r_1>r_2$ when $D$ is a similarity measure.

We will briefly describe, from \cite{IM}, how a LSH family can be used to
solve the $(R,c)$-NN problem: We choose $r_1=R$ and $r_2=c \cdot R$.
Given a family ${\cal H}$ of hash functions with parameters
$(r_1,r_2,p_1,p_2)$ as in Definition~\ref{defn:lsh}, we
amplify the gap between the ``high'' probability $p_1$ and ``low''
probability $p_2$ by concatenating several functions.  In particular, for
$k$ specified later, define a function family ${\cal G} = \{g:S \to U^k\}$
such that $g(v)=(h_1(v), \ldots ,h_k(v))$, where $h_i \in {\cal H}$. For
an integer $L$ we choose $L$ functions $g_1, \ldots, g_L$ from ${\cal G}$,
independently and uniformly at random.  During preprocessing, we store
each $v \in P$ (input point set) in the bucket $g_j(v)$, for $j=1, \ldots,
L$. Since the total number of buckets may be large, we retain only the
non-empty buckets by resorting to hashing.
%If any bucket
%contains more than one element, we retain an arbitrary one.
To process a query $q$, we search all buckets $g_1(q), \ldots, g_L(q)$; as
it is possible (though unlikely) that the total number of points stored in
those buckets is large, we interrupt search after finding first $3L$ points
(including duplicates).  Let $v_1, \ldots, v_t$ be the points encountered
therein.
%For each $p_j$, we check whether $p_j \in B(q,r_2)$ and if that is the case,
%we return {\sc yes} (together with $p_j$); if no such $p_j$ is found, we
%return {\sc no}.
For each $v_j$, if $v_j \in B(q,r_2)$ then we return {\sc yes} and $v_j$,
else we return {\sc no}.

%Let $W_b(q)=P - B(q,b)$, and $p^*$ be the point in $P$ closest to $q$.
The parameters $k$ and $L$ are chosen so as to ensure that with a constant probability the following two properties hold:
%there exists $g_j$ such that the following properties hold:
\begin{enumerate}
\item If there exists $v^* \in B(q,r_1)$ then $g_j(v^*)=g_j(q)$ for some $j=1 \ldots L$, and
\item The total number of collisions of $q$ with points from $P - B(q,r_2)$ is less than $3L$, i.e.
\[ \sum_{j=1}^L |(P-B(q,r_2)) \cap g_j^{-1}(g_j(q))| < 3L.\]
%$g_j(p') \neq g_j(q)$, for all $p' \in P - B(q,r_2)$ and $j=1 \ldots
\end{enumerate}

Observe that if (1) and (2) hold, then the algorithm is correct.
%{\bf Piotr: Verify this part}
It follows (see \cite{IM} Theorem~5 for details) that if we set $k =
\log_{1/p_2} n$, and $L = n^\rho$ where $\rho=\frac{\ln 1/p_1}{\ln 1/p_2}$
then (1) and (2) hold with a constant probability.
\iffalse
 In order to boost this
probability and get a high probability result we run $O(\log n)$ parallel
instances of the algorithm. Moreover, each hash function evaluation
requires $O(d)$ time.
\fi
Thus, we get following theorem (slightly different
version of Theorem~5 in \cite{IM}), which relates the efficiency of
solving $(R,c)$-NN problem to the sensitivity parameters of the LSH.
\begin{theorem}
\label{t:sens}
Suppose there is a $(R,c R,p_1,p_2)$-sensitive family ${\cal H}$ for a
distance measure $D$.
Then there exists an algorithm for $(R,c)$-NN under measure $D$
which uses  $O(d n +  n^{1+\rho})$ space, with query time dominated by
$O(n^{\rho})$ distance computations, and $O(n^{\rho} \log_{1/p_2} n)$ evaluations of
hash functions from ${\cal H}$,
%evaluations of the hash function for each query,
where $\rho=\frac{\ln 1/p_1}{\ln 1/p_2}$.
\end{theorem}

% besides
% removing the polylog factors from the query time it strictly improves the
% exponent of $n$ in space requirement and query time.


\iffalse
\section{Preliminaries}
Dynamic $c$-approximate {\em near neighbor} problem: for a set of points $P$,
construct a data structure that given any query point $q$:
\begin{itemize}
\item If there is a point $s$ within a predefined
distance $R$ from $q$, report a (possibly different) point $p'$ within distance $cR$ from $q$
\item If no such point exists, any answer is fine
\end{itemize}

Since we consider points in $l_p^d$, without loss of generality we can consider $R=1$,
which we assume from now on.
\fi

\section{Our LSH scheme}


 In this section, we present a LSH family based on
$p$-stable distributions, that works for all $p \in (0,2]$. 

Since we consider points in
$l_p^d$, without loss of generality we can consider $R=1$, which we assume
from now on.

\iffalse
We prove the
following main theorem about the sensitivity parameters of this hash
family:

\begin{theorem}
\label{thm:general_p}
For any $p \in (0,2]$ there is a $(r_1,r_2,p_1,p_2)$-sensitive family
${\cal H}$ for $l_p^d$ such that for any $\gamma > 0$,
$$\rho=\frac{\ln 1/p_1}{\ln 1/p_2} \le (1+\gamma)\cdot\max\left(\frac{1}{c^p}, \frac{1}{c}\right).$$
\end{theorem}

Thus, using this hashing scheme we get a data structure for $(R,c)$-NN
problem in $l_p^d$ whose space and query time has exponent $\rho$ which is
arbitrarily close to $\max\left(\frac{1}{c^p}, \frac{1}{c}\right)$. 
\fi


\subsection{$p$-stable distributions}

Stable distributions \cite{Zol} are defined as limits of normalized sums
of independent identically distributed variables (an alternate definition
follows).  The most well-known example of a stable distribution is
Gaussian (or normal) distribution.  However, the class is much wider; for
example, it includes heavy-tailed distributions.

{\noindent \bf Stable Distribution:} A distribution ${\cal D}$ over $\Re$ is called {\em $p$-stable}, if there exists $p \ge 0$ such that for any $n$ real numbers $v_1 \ldots v_n$ and i.i.d. variables $X_1 \ldots X_n$ with distribution ${\cal D}$, the random variable $\sum_i v_i X_i$ has the same distribution as the variable $(\sum_i |v_i|^p)^{1/p} X$, where $X$ is a random variable with distribution ${\cal D}$.

It is known~\cite{Zol} that stable distributions exist for any $p \in (0,2]$.
In particular:
\begin{itemize}
\item a {\em Cauchy distribution} ${\cal D}_C$, defined by the density function
$c(x)=\frac{1}{\pi} \frac{1}{1+x^2}$, is $1$-stable
\item a {\em Gaussian (normal) distribution} ${\cal D}_G$, defined by the density function $g(x)= \frac{1}{\sqrt{2\pi}} e^{-x^2/2}$, is $2$-stable
\end{itemize}

\iffalse
While there are no closed form expressions for the density functions $f_p$
and/or distribution functions $F_p$ of general $p$-stable distributions,
there is an estimation for these functions known as the Pareto
estimation~\cite{Nol}.  As per this estimation:

\begin{tabular}{c}
\hfill~\\
$\forall\delta>0\ \exists x_0\mathrm{\ s.t.\ }\forall x\geq x_0,$\hfill~\\
~~$C_px^{-p}(1-\delta)\leq F_p(x)\leq C_px^{-p}(1+\delta)$\hfill~\\
\hfill~\\
$\forall\delta>0\ \exists x_0\mathrm{\ s.t.\ }\forall x\geq x_0,$\hfill~\\
~~$C_px^{-(p+1)}(1-\delta)\leq f_p(x)\leq C_px^{-(p+1)}(1+\delta)$\hfill~\\
\hfill~\\
\end{tabular}
%$$\forall\delta>0\ \exists x_0\mathrm{\ s.t.\ }\forall x\geq x_0,
%C_px^{-p}(1-\delta)\leq F_p(x)\leq C_px^{-p}(1+\delta)$$
%$$\forall\delta'>0\ \exists x_0'\mathrm{\ s.t.\ }\forall x\geq x_0',
%C_px^{-(p+1)}(1-\delta')\leq f_p(x)\leq C_px^{-(p+1)}(1+\delta')$$ 

\noindent
where $C_p=\frac{1}{\pi}\Gamma(p) \sin(\pi p/2)$.  From this
estimation, it is clear that the expectation ($\int_0^{\infty} t f_p(t)dt$)
is finite for $p>1$ and
infinite for $p<1$.  These facts will become important later in our
analysis.  Moreover, this implies that $p$-stable distributions do not
have finite second moments or finite variances (except for $p=2$).
\fi

We note from a practical point of view, despite the lack of closed form
density and distribution functions, it is known~\cite{CMS} that one can
generate $p$-stable random variables essentially from two independent
variables distributed uniformly over $[0,1]$.

Stable distribution have
found numerous applications in various fields (see the survey~\cite{Nol}
for more details). In computer science, stable distributions were used for
``sketching'' of high dimensional vectors by Indyk (\cite{I00b}) and since
have found use in various applications.
The main property of $p$-stable distributions mentioned in the definition
above directly translates into a
sketching technique for high dimensional vectors. The idea is to
generate a random vector $\bm{a}$ of dimension $d$ whose each entry is
chosen independently from a $p$-stable distribution. Given a vector
$\bm{v}$ of dimension $d$, the dot product $\bm{a}.\bm{v}$ is a random variable
which is distributed as $(\sum_i |v_i|^p)^{1/p} X$ (i.e., $||\bm{v}||_p
X$), where $X$ is a random variable with $p$-stable distribution. A small
collection of such dot products ($\bm{a}.\bm{v}$), corresponding to different
$\bm{a}$'s, is termed as the sketch of the vector $\bm{v}$ and can be used to
estimate $||\bm{v}||_p$ (see
\cite{I00b} for details). It is easy to see that such a sketch is linearly
composable, i.e. $\bm{a}.(\bm{v_1} -\bm{v_2}) = \bm{a}.\bm{v_1} - \bm{a}.\bm{v_2}$.

\subsection{Hash family}
\label{subsec:family}
In this paper we use $p$-stable distributions in a slightly different
manner. Instead of using the dot products ($\bm{a}.\bm{v}$) to estimate
the $l_p$ norm we use them to assign a hash value to each vector
$\bm{v}$. Intuitively, the hash function family should be locality
sensitive, i.e. if two vectors ($\bm{v_1}, \bm{v_2}$) are close (small
$||\bm{v_1}-\bm{v_2}||_p$) then they should collide (hash to the same
value) with high probability and if they are far they should collide with
small probability. The dot product $\bm{a}.\bm{v}$ projects each vector to the
real line; It follows from $p$-stability that for two vectors
($\bm{v_1},\bm{v_2}$) the distance between their projections
($\bm{a}.\bm{v_1} - \bm{a}.\bm{v_2}$) is distributed as $||\bm{v_1}-\bm{v_2}||_p X$
where $X$ is a $p$-stable distribution. If we ``chop'' the real line into
equi-width segments of appropriate size $r$ and assign hash values to vectors
based on which segment they project onto, then it is intuitively clear
that this hash function will be locality preserving in the sense described
above.

Formally, each hash function $h_{\bm{a},b}(\bm{v}): {\cal R}^d \rightarrow
{\cal N}$ maps a $d$ dimensional vector $\bm{v}$ onto the set of
integers. Each hash function in the family is indexed by a choice of
random $\bm{a}$ and $b$ where $\bm{a}$ is, as before, a $d$ dimensional
vector with entries chosen independently from a $p$-stable
distribution and $b$ is a real number chosen uniformly from the range
$[0,r]$. For a fixed $\bm{a},b$ the hash function $h_{\bm{a},b}$ is given by
$h_{\bm{a},b}(\bm{v})=\lfloor{\bm{a} \cdot \bm{v}+b\over r}\rfloor$

Next, we compute the probability that two vectors $\bm{v_1}, \bm{v_2}$
collide under a hash function drawn uniformly at random from this
family. Let $f_p(t)$ denote the probability density function of the
{\bf absolute value} of the $p$-stable distribution. We may drop the subscript
$p$ whenever it is clear from the context.
\iffalse
 For example, the density
functions for $p=1$ (Cauchy) and $p=2$ (Gaussian) are $f_1(t) =
\frac{2}{\pi} \frac{1}{1+t^2}$ and $f_2(t) = \frac{2}{\sqrt{2\pi}}
e^{-t^2/2}$ respectively. For a general $p$, there may not exist such a
closed form expression for the density function.
\fi
 For the two vectors $\bm{v_1},
\bm{v_2}$, let $c = ||\bm{v_1} - \bm{v_2}||_p$. For a random vector $\bm{a}$ whose entries are
drawn from a $p$-stable distribution, $\bm{a}.\bm{v_1} -
\bm{a}.\bm{v_2}$ is distributed as $c X$ where $X$ is a random variable
drawn from a $p$-stable distribution. Since $b$ is
drawn uniformly from $[0,r]$ it is easy to see that

%\begin{eqnarray*}
%\label{eq:prob}
$$p(c)=Pr_{\bm{a},b} [h_{\bm{a},b}(\bm{v_1}) = h_{\bm{a},b}(\bm{v_2})] =
%Pr_{\bm{a}}[|\bm{a}\cdot (\bm{v_1} - \bm{v_2})| =t]
%Pr_b[h_{\bm{a},b}(\bm{v_1}) = h_{\bm{a},b}(\bm{v_2}) \mid |\bm{a}\cdot (\bm{v_1} - \bm{v_2})| =
%t]
%\cr &=&
\int_0^r \frac{1}{c}
f_p(\frac{t}{c}) (1-\frac{t}{r}) dt$$
%\end{eqnarray*}

For a fixed parameter $r$ the probability of collision decreases
monotonically with $c=||\bm{v_1} - \bm{v_2}||_p$. Thus, as per
Definition~\ref{defn:lsh} the family of hash functions above is $(r_1,
r_2, p_1, p_2)$-sensitive for $p_1 = p(1)$ 
%\int_0^r f_p(t) (1-\frac{t}{r}) dt$ 
and $p_2 = p(c)$
%\int_0^r \frac{1}{c} f_p(\frac{t}{c}) (1-\frac{t}{r}) dt$ 
for $r_2/r_1 = c$.

In what follows we will bound the ratio $\rho=\frac{\ln 1/p_1}{\ln
1/p_2}$, which as discussed earlier is critical to the performance when
this hash family is used to solve the $(R,c)$-NN problem.

Note that we have not specified the parameter $r$, for it depends
on the value of $c$ and $p$. For every $c$ we would like to choose a finite $r$ that
makes $\rho$ as small as possible.
\iffalse
Note also, the way we have described the hash family it
appears that we may require a large number (potentially infinite) of random
bits to represent each function from this family: Each hash function is
specified by a random vector $\bm{a}$ and real $b$. {\bf Piotr : Verify this}
However, using techniques similar to those in Indyk \cite{I00b} and
Engebretsen et al~\cite{EIO}, we can reduce the number of random bits
required to $O(\log^2 n)$, where $n$ is the number of vectors in the data
set. The techniques involve the use of Nisan's pseudorandom number
generator for space bounded computation~\cite{Nis}.
\fi

\iffalse
We want to solve the Near Neighbor problem for $l_p$ norms,
$p\in(0,2]$, using the locality sensitive hashing
approach.
We use hash function $h_{a,b}(s)=\lfloor{a \cdot s+b\over r}\rfloor$ where $a$ is a
$d$-dimensional vector drawn from a $p$-stable distribution and
$b$ is chosen uniformly at random from $[0,r]$.

\begin{defn}
Let $E_c$ be the event that two close points $v_1$ and $v_2$ collide,
i.e. $[h_{a,b}(v_1)=h_{a,b}(v_2)  : ||v_1-v_2||_p < 1],$ and $p_1=\Pr_{a,b}[E_c]$.
Let $E_f$ be the event that two far points $v_1'$ and $v_2'$ collide,
i.e. $[h_{a,b}(v_1') = h_{a,b}(v_2') : ||v_1'-v_2'||_p > c],$ and $p_2=\Pr_a[E_f]$.
\end{defn}

Now conditioning on the distance of points $s$ and $q$ in the projection
onto line $a$, we see
\begin{eqnarray}
p_1
&\equiv&
\Pr_a[E_c]
\cr&=&
\int_0^r\Pr_a[E_c : |a\cdot(q-s)|=t]\Pr_a[|a\cdot(q-s)|=t]dt
\cr&=&
\int_0^r(1-\frac{t}{r})\Pr_a[|a\cdot(q-s)|=t]dt
\end{eqnarray}
As $a$ is drawn from a $p$-stable distribution, it is equivalent, in
distribution, to the random variable $||s-q||_px$ where $x$ is again
drawn according to a $p$-stable distribution.  Let $f(\cdot)$ be the
density function of the absolute value of a $p$-stable random variable.
Using the assumption
that $||s-q||_p \le 1$
% and the symmetry of the $p$-stable distributions around $x=0$,
we see
$$p_1>\int_0^r(1-\frac{t}{r})f(\frac{t}{1})dt.$$
Similarly,
$$p_2<\int_0^r(1-\frac{t}{r})\frac{1}{c}f(\frac{t}{c}) dt.$$
\fi

\section{Computational analysis of the ratio $\rho=\frac{\ln 1/p_1}{\ln 1/p_2}$}
\label{s:comp}

In this section we focus on the cases of $p=1,2$. In these cases the 
ratio $\rho$ can be explicitly evaluated.
We compute and plot this ratio
and compare it with $1/c$. Note, $1/c$ is the best (smallest) known
exponent for $n$ in the space requirement and query time that is achieved
in \cite{IM} for these cases. 
\iffalse
For the case of $p=1/2$, we evaluate $\rho$
numerically and compare it with $1/\sqrt{c}$ for various values of $r$
and $c$.  
For the general case we prove that there
exists a value of $r$ for which the ratio $\rho$ is arbitrarily close to
$\max\left(\frac{1}{c^p},
\frac{1}{c}\right)$, thereby proving Theorem~\ref{thm:general_p}
\fi

\subsection{Computing the ratio $\rho$ for special cases} For the special
cases $p=1,2$ we can compute the probabilities $p_1,p_2$, using the
density functions mentioned before. A simple
calculation shows that $p_2 = 2\frac{tan^{-1}({r/c})}{\pi} - \frac{1}{\pi
(r/c)} \ln ({1 + (r/c)^2})$ for $p=1$ (Cauchy) and
$p_2 = 1 - 2 norm(-r/c) -
\frac{2}{\sqrt{2\pi}r/c}(1 - e^{-(r^2/2c^2)})$ for $p=2$ (Gaussian), where
$norm(\cdot)$ is the cumulative distribution function (cdf) for
a random variable that is distributed as $N(0,1)$.
The value of $p_1$ can be obtained
by substituting $c=1$ in the formulas above.

\begin{figure*}[t]
\begin{center}
\mbox{
subfigure[Optimal $\rho$ for $l_1$]{
IncludeGraphics[width=7.5cm,height=5.6cm]{l1 plot.eps}
\label{fig:l1_plot}
} \quad
subfigure Optimal $\rho$ for $l_2$]{
includegrphics[width=7.5cm,height=.6cm]{l2 plot.eps}
\label{fig:l2_plot}
}
}
\caption{Optimal $\rho$ vs $c$}
\label{fig:opt_rho}
\end{center}
\end{figure*}

For $c$ values in the range $[1,10]$ (in increments of $0.05$) we compute
the minimum value of $\rho$, $\rho(c) = min_r \log ({1/p_1})/\log
({1/p_2})$, using {\em Matlab}. The plot of $c$ versus $\rho(c)$ is shown
in Figure~\ref{fig:opt_rho}. The crucial observation for the case $p=2$ is
that the curve corresponding to optimal ratio $\rho$ ($\rho(c)$) lies
strictly below the curve $1/c$. As mentioned earlier, this is a strict
improvement over the previous best known exponent $1/c$ from
\cite{IM}. While we have computed here $\rho(c)$ for $c$
in the range $[1,10]$, we believe that $\rho(c)$ is strictly less than
$1/c$ for all values of $c$. 
%Proving this is part of our plans for future work. 

For the case $p=1$, we observe that $\rho(c)$ curve is very close to
$1/c$, although it lies above it. The optimal $\rho(c)$ was
computed using {\em Matlab} as mentioned before. The {\em Matlab} program
has a limit on the number of iterations it performs to compute the minimum
of a function. We reached this limit during the computations. If we
compute the true minimum, then we suspect that it will be very close to
$1/c$, possibly equal to $1/c$, and that this minimum might be reached at
$r = \infty$.

\begin{figure*}[t]
\begin{center}
\mbox{
subfigure[$\rho$ vs $r$ for $l_1$]{
includegrphics[width=5.6cm,height=7.5cm,angle=-90]{$r_var_l1.ps$}
\label{Fig:r_var_l1}
} \quad
subfigure[$\rho$ vs $r$ for $l_2$]{
includegrphics[width=5.6cm,height=7.5cm,angle=-90]{$r_var_l2.ps$}
\label{Fig:r_var_l2}
}
}
\caption{$\rho$ vs $r$}
\label{fig:r_var}
\end{center}
\end{figure*}

If one were to implement our LSH scheme, ideally they would want to know
the optimal value of $r$ for every $c$.  For $p=2$, for a given value
of $c$,
% unless $c$ is very large (and hence not of practical interest),
we can compute the value of $r$ that gives the optimal value of
$\rho(c)$. This can be done using programs like {\em Matlab}.
However, we
observe that for a fixed $c$ the value of $\rho$ as a function of $r$
is more or less stable after a certain point
(see Figure~\ref{fig:r_var}).
%Figure~\ref{fig:r_var} shows the value of $\rho$
%as function of $r$ for some fixed values of $c$.
Thus, we observe that
$\rho$ is not very sensitive to $r$ beyond a certain point and as
long we choose $r$ ``sufficiently'' away from $0$, the $\rho$ value will be
close to optimal. Note, however that we should not choose an $r$ value that
is too large. As $r$ increases, both $p_1$ and $p_2$ get closer to
$1$. This increases the query time, since $k$, which is the
``width'' of each hash function (refer to Subsection~\ref{subsec:lsh}),
increases as $\log_{1/p_2} n$.

We mention that for the $l_2$ norm, the optimal value of $r$ appears to be a
(finite) function of $c$.

\begin{figure*}[t]
\begin{center}
\mbox{
subfigure[$\rho$ vs $c$ for $l_1$]{
includegrphics[width=5.6cm,height=7.5cm,angle=-90]{$c_var_l1.ps$}
\label{Fig:c_var_l1}
} \quad
subfigure[$\rho$ vs $c$ for $l_2$]{
includegrphics[width=5.6cm,height=7.5cm,angle=-90]{$c_var_l2.ps$}
\label{Fig:c_var_l2}
}
}
\caption{$\rho$ vs $c$}
\label{fig:c_var}
\end{center}
\end{figure*}

We also plot $\rho$ as a function of $c$ for a few fixed $r$ values(See
Figure~\ref{fig:c_var}). For $p=2$, we observe that for moderate $r$
values the $\rho$ curve ``beats'' the $1/c$ curve over a large range of
$c$ that is of practical interest. For $p=1$, we observe that as $r$
increases the $\rho$ curve drops lower and gets closer and closer to the
$1/c$ curve.

\iffalse
For general $p$, we can numerically calculate the density function from
existing software~\cite{Nol1}.  Then, using Maple, we can compute $\rho$
%the exponent in the performance of the algorithm based on
for various settings of the parameters $c$ and $r$.  

We provide two plots
for the case $p=0.5$ (Figure~\ref{fig:levy}). In
Figure~\ref{fig:r_p0.500} we plot $\rho$ as a function of $c$ for a few
fixed values of $r$ and compare it to $1/\sqrt{c}$.  Notice for $r=1$,
$\rho$ is worse than $1/\sqrt{c}$, (i.e. $\gamma>0$ in
Theorem~\ref{thm:general_p}).  However, for $r=25$, $\rho$ is about
$1/\sqrt{c}$, and if we allow $r=75$, we can beat $1/\sqrt{c}$ for this
range of $c$ (for larger $c$, this might not be case).  In fact,
for a fixed value of $c$ ($c=2$) we observe the tradeoff in
performance as we increase $r$.  Figure~\ref{fig:c_p0.500} shows $\rho$ as
a function of $r$ for $c=2$ and compares it to $1/\sqrt{c}$.  For this
value of $c$, we note the performance matches $1/\sqrt{c}$ for relatively
small $r\simeq25$.  As the plot shows, we can decrease $\rho$ by
increasing $r$.
\fi
\iffalse
 It is important to note the
constants hidden in the big-$O$ notation of Theorem~\ref{t:sens} grow as
$r$ grows, and so our performance can not be arbitrarily good.  This
suggests that there is some optimal $r$ which appropriately balances the
constants and the $n^\rho$ terms in the query time.  Unfortunately, we
can not calculate this for general $p$ as we do not know the density
function and thus do not know how the constants depend on $r$.
\fi

\iffalse
\begin{figure*}[t]
\begin{center}
\mbox{
subfigure[$\rho$ vs $c$ for $l_{1/2}$]{
includegrphics[width=5.6cm,height=7.5cm]{p05001.eps}
\label{fig:r_p0.500}
} \quad
subfigure[$\rho$ vs $r$ for $l_{1/2}$]{
includegrphics[width=5.6cm,height=7.5cm]{p05002.eps}
\label{fig:c_p0.500}
}
}
\caption{$\rho$ for $l_{1/2}$}
\label{fig:levy}
\end{center}
\end{figure*}
\fi

\section{Empirical Evaluation of our Technique}
\label{Sec:expt} 
\iffalse
We feel that our locality sensitive hashing based
technique for near neighbor queries will be an attractive choice
for implementing a nearest neighbor data structures. The two main
advantages of our technique, over all other previously known
techniques, are its generality and simplicity. As mentioned
earlier, our technique is the only one that is known to work for
high-dimensional nearest neighbor problem over $l_p$ norms with $p
< 1$. Similarity search under such {\em fractional} norms have
recently attracted interest~\cite{AHK,CIKM}.
\fi

In this section we present an experimental evaluation of our novel LSH
scheme.  We focus on the Euclidean norm case, since this occurs most
frequently in practice.  Our data structure is implemented for main
memory.  
%In our experiments we observed that our technique is an order
%of magnitude faster than the previously best known main memory
%implementation for nearest neighbor search based on the $kd$-tree
%data structure~\cite{AMexp}.

In what follows, we briefly discuss some of the issues pertaining to the
implementation of our technique.  We then report some preliminary
performance results based on an empirical comparison of our technique to
the $kd$-tree data structure.

{\noindent \bf Parameters and Performance Tradeoffs:} The three
main parameters that affect the performance of our algorithm are:
number of projections per hash value ($k$), number of hash tables
($l$) and the width of the projection ($r$).
In general, one could also introduce another parameter (say $T$),
such that the query procedure stops after retrieving $T$ points.
In our analysis, $T$ was set to $3l$. 
In our experiments, however, the query procedure retrieved {\em all} 
points colliding with the query (i.e., we used $T=\infty$).
This reduces the number of parameters and simplifies the choice of the
optimal.

\iffalse
%({\bf I am using 'r' and
%not 'w' to be consistent with SOCG write up. If you change to 'w'
%make sure to change that as well}).
 As we increase the value of
$k$ the probability of two points colliding decreases
exponentially. This decreases the number of false positives per
hash tables, i.e. points that are not near neighbors but hash into
the same bucket as the query point. At the same time, increasing
$k$ also increases the number of false negatives, i.e. points that
are true near neighbors but do not hash into the same bucket as
the query point. If we increase $k$, we also have to increase $l$
in order to maintain the fraction of false negatives below a
certain threshold. As a result, not only does the time to compute
a single hash function increase, but we also need to evaluate more
hash functions per query point. However, the total number of data
points that collide with the query point over {\bf all} the hash
functions, i.e. candidate near neighbors, is reduced. Effectively,
we spend more time doing hashing and less time checking distances
with the candidate near neighbors. Thus, in applications where it
is costly to do distance computations, our algorithm provides a
tradeoff of increased hashing time for decreased distance
computation time. Such a scenario can occur if the data points
reside on disk or if the dimensionality of the data set is too
high. In contrast, data structures like $R$-trees or
$kd$-trees do not provide such tradeoff and spend all their
time doing distance computations with data points.
\fi

For a given value of $k$, it is easy to find the optimal value of $l$
which will guarantee that the fraction of false negatives are no more than
a user specified threshold. This process is exactly the same as in an
earlier paper by Cohen et al. (\cite{eight}) that uses locality sensitive
hashing to find similar column pairs in market-basket data, with the
similarity exceeding a certain user specified threshold. In our
experiments we tried a few values of $k$ (between $1$ and $10$) and below
we report the $k$ that gives the best tradeoff for our scenario. The
parameter $k$ represents a tradeoff between the time spent in computing
hash values and time spent in pruning false positives, i.e. computing
distances between the query and candidates; a bigger $k$ value 
increases the number of hash computations. In general we could do a binary
search over a large range to find the optimal $k$ value. This binary
search can be avoided if we have a good model of the relative times of
hash computations to distance computations for the application at hand.

Decreasing the width of the projection ($r$) decreases the probability
of collision for any two points. Thus, it has the same effect as
increasing $k$. As a result, we would like to set $r$ as small as
possible and in this way decrease the number of projections we need to
make.  However, decreasing $r$ below a certain threshold increases the
quantity $\rho$, thereby requiring us to increase $l$. Thus we cannot
decrease $r$ by too much. For the $l_2$ norm we found the optimal value
of $r$ using Matlab which we used in our experiments.

\iffalse
With the right set of parameters, as described above, we found
that the query processing time for our algorithm is an order of
magnitude faster than the best main memory implementation of
$kd$-tree (available at {\tt http://www.cs.umd.edu/~mount/ANN/})
against which we compared it.
 Moreover, the advantage of our
algorithm increases by a lot with increasing number of data points
and dimensionality. Thus it scales much better than $kd$-tree. We
feel that our algorithm will be an attractive choice of
implementation for massive and high dimensional data sets. 
\fi
Before
we report our performance numbers we will next describe the data
set and query set that we used for testing.

{\noindent \bf Data Set:} We used synthetically generated data sets and
query points to test our algorithm. The dimensionality of the underlying
$l_2$ space was varied between $20$ and $500$. We considered generating
all the data and query points independently at random. Thus, for a data
point (or query point) its coordinate along every dimension would be
chosen independently and uniformly at random from a certain range $[-a,
a]$. However, if we did that, given a query point all the data points
would be sharply concentrated at the same distance from the query point
as we are operating in high dimensions. Therefore, approximate
nearest neighbor search would not make sense on such a data set. Testing
approximate nearest neighbor requires that for every query point $q$,
there are few data points within distance $R$ from $q$ and most of the
points are at a distance no less than $(1+ \epsilon) R$.  We call this a
``planted nearest neighbor model''.


In order to ensure this property we generate our points as follows 
(a similar approach was used in~\cite{Yia2}). We
first generate the query points at random, as above. We then generate
the data points in such a way that for every query point, we guarantee
at least a single point within distance $R$ and all other points are
distance no less than $(1+ \epsilon) R$.  This novel way of
generating data sets ensures every query point has a
%is appropriate for testing near neighbor data structures. 
few (in our case, just one) approximate nearest neighbors, while most
points are far from the query.

The resulting data set has several interesting properties. Firstly, it 
constitutes the worst-case input to LSH (since there is only one 
correct nearest neighbor, and all other points are ``almost'' correct
nearest neighbors). Moreover, it captures the typical situation occurring in real life similarity search applications, in which there are few points that are relatively close to the query point, and most of the database points lie quite far from the query point.



For our experiments the range $[-a, a]$ was set to $[-50, 50]$.
%so that each dimension could be stored as a char. 
The total number of data points was varied between $10^4$ and
$10^5$. Both our algorithm and the $kd$-tree take as input the
approximation factor $c = (1+ \epsilon)$. However, in addition to $c$
our algorithm also requires as input the value of the distance $R$
(upper bound) to the nearest neighbor. This can be avoided by guessing
the value of $R$ and doing a binary search.  We feel that for most real
life applications it is easy to guess a range for $R$ that is not too
large. As a result the additional multiplicative overhead of doing a
binary search should not be much and will not cancel the gains that we
report.

{\bf Experimental Results:} We did three sets of experiments to evaluate
the performance of our algorithm versus that of $kd$-tree: we
increased the number $n$ of data points, the dimensionality $d$ of the data
set, and the approximation factor $c=(1+ \epsilon)$. In each set of
experiments we report the average query processing times for our
algorithm and the $kd$-tree algorithm, and also the ratio of the two
((average query time for $kd$-tree)/( average query time for our
algorithm)), i.e. the speedup achieved by our algorithm.  We ran our
experiments on a Sun workstation with 650 MHz UltraSPARC-IIi, 512KB L2
cache processor, having no special support for vector computations, with
512 MB of main memory.


 For all our experiments we set the parameters $k=10$ and $\ell =
30$. Moreover, we set the percentage of false negatives that we can
tolerate up to $10\%$ and indeed for all the experiments that we report
below we did not get the more than $7.5\%$ false negatives, in fact less
in most cases.



\begin{figure*}[t]
\begin{center}
\mbox{
subfigure[query time vs $n$]{
%includegrphics[width=7.5cm,height=5.6cm]{n_time.eps}
\label{Fig:n-speedup1}
} \quad
subfigure[speedup vs $n$]{
%includegrphics[width=7.5cm,height=5.6cm]{$n_speedup.eps$}
\label{Fig:n-speedup2}
}
}
\caption{Gain as data size varies}
\label{Fig:n-speedup}
\end{center}
\end{figure*}

For all the query time graphs that we present, the curve that lies above
is that of $kd$-tree and the one below is for our algorithm.

For the first experiment we fixed $\epsilon = 1$, $d = 100$
and $r = 4$ (the width of projection).  We varied the number of data points from $10^4$ to
$10^5$. Figures~\ref{Fig:n-speedup1}~and~\ref{Fig:n-speedup2} show
the processing times and speedup respectively as $n$ is varied. As
we see from the Figures, the speedup seems to increase linearly
with $n$.
% Thus, our algorithm scales much better than the
%$kd$-tree algorithm.

\begin{figure*}[t]
\begin{center}
\mbox{
subfigure[query time vs dimension]{
includegrphics[width=7.5cm,height=5.6cm]{$dim_time.eps$}
\label{Fig:dim-speedup1}
} \quad
subfigure[speedup vs dimension]{
includegrphics[width=7.5cm,height=5.6cm]{$dim_speedup.eps$}
\label{Fig:dim-speedup2}
}
}
\caption{Gain as dimension varies}
\label{Fig:dim-speedup}
\end{center}
\end{figure*}

For the second experiment we fixed $\epsilon = 1$, $n = 10^5$ and
$r = 4$. We varied the dimensionality of the data set from $20$ to
$500$. Figures~\ref{Fig:dim-speedup1}~and~\ref{Fig:dim-speedup2}
show the processing times and speedup respectively as $d$ is
varied. As we see from the Figures, the speedup seems to increase
with the dimension. 
%Again, this suggests that out algorithm scales
%better than $kd$-tree with the number of dimensions.

\begin{figure*}[t]
\begin{center}
\mbox{
subfigure[query time vs $\epsilon$]{
includegrphics[width=7.5cm,height=5.6cm]{$eps_time.eps$}
\label{Fig:eps-speedup1}
} \quad
subfigure[speedup vs $\epsilon$]{
includegrphics[width=7.5cm,height=5.6cm]{$eps_speedup.eps$}
\label{Fig:eps-speedup2}
}
}
\caption{Gain as $\epsilon$ varies}
\label{Fig:eps-speedup}
\end{center}
\end{figure*}

For the third experiment we fixed $n = 10^5$ and $d = 100$. The
approximation factor $(1+\epsilon)$ was varied from $1.5$ to $4$.
The width $r$ was set appropriately as a function of $\epsilon$.
Figures~\ref{Fig:eps-speedup1}~and~\ref{Fig:eps-speedup2} show the
processing times and speedup respectively as $\epsilon$ is varied.
%As we see from the Figures, the speedup mostly increases with $\epsilon$.
%({\bf this is incorrect as of now. I am hoping it will change}).

{\noindent \bf Memory Requirement:} The memory requirement for our
algorithm equals the memory to store the data points themselves
and the memory required to store the hash tables. From our
experiments, typical values of $k$ and $l$ 
are $10$ and $30$ respectively. If we insert each point
in the hash tables along with their hash values and a pointer to
the data point itself, it will require $l\cdot(k+1)$ words (int)
of memory, which for our typical $k,l$ values evaluates to $330$
words. We can reduce the memory requirement by not storing the
hash value explicitly as concatenation of $k$ projections, but
instead hash these $k$ values in turn to get a single word for the
hash. This would reduce the memory requirement to $l\cdot 2$, i.e.
$60$ words per data point. If the data points belong to a high
dimensional space (e.g., with $500$ dimension or more), then
the overhead of maintaining the hash table is not much (around
$12\%$ with the optimization above) as compared to storing the
points themselves. Thus, the memory overhead of our algorithm is
small. 
%Figure~\ref{Fig:mem} shows the memory requirement of our
%algorithm as $n$ is varied. ({\bf I dont know if we can get such a
%figure and if we should since our numbers dont seem to conform
%with above discussion. Take a decision and edit accordingly}).

%({\bf I feel that the discussion about ``stop when 3L'' should be
%in the section that describes the algorithm.})

\section{Conclusions}

In this paper we present a new LSH scheme for the similarity search in high-dimensional spaces.
The algorithm is easy to implement, and generalizes to arbitrary $l_p$ norm, for $p \in [0,2]$.
We provide theoretical, computational and experimental evaluations of the algorithm.

Although the experimental comparison of LSH and kd-tree-based algorithm suggests that the former outperforms the latter, there are several caveats that one needs to keep in mind:
\begin{itemize}
\item We used the kd-tree structure ``as is''. Tweaking its parameters would likely improve its performance.
\item LSH solves the decision version of the nearest neighbor problem, while kd-tree solves the optimization version.
Although the latter reduces to the former, the reduction overhead increases the running time.
\item One could run the approximate kd-tree algorithm with approximation parameter $c$ that is much larger than the intended approximation.
Although the resulting algorithm would provide very weak guarantee on the quality of the returned neighbor, typically the actual error is much smaller than the guarantee.
\end{itemize}



\bibliographystyle{plain}

{\footnotesize

\bibliography{paper.bib}

}



















\appendix


\section{Growth-restricted data sets}

In this section we focus exclusively on data sets living in the $l_2^d$ norm.

Consider a data set $P$, query $q$, and let $p$ be the closest point in $P$ to $q$.
Assume we know the distance $\|p-q\|$, in which case we can assume that it is equal to $1$, by scaling\footnote{Similar guarantees can be proved when we only know a constant approximation to the distance.}.
For $c \ge 1$, let $P(q,c)=P \cap B(q,c)$ and let $N(q,c)=|P(q,c)|$.


We consider a ``single shot'' LSH algorithm, i.e., one that uses only $L=1$ indices, but examines all points in the bucket containing $q$.
We use the parameters $k=r= T \log n$, for some constant $T>1$.
This implies that the hash function can be evaluated in time $O(\log n)$.
%Combined with the theorem below, this gives an algorithm with query time $O(d \log n + d 2^{O(b)})$.

\begin{theorem}
If $N(q,c)=O(c^b)$ for some $b>1$, then the ``single shot'' LSH algorithm finds $p$ with constant probability in expected time $d (\log n + 2^{O(b)})$. 
\end{theorem}

\noindent{\bf Proof:} For any point $p'$ such that $\|p'-q\|=c$, 
the probability that $h(p')=h(q)$ is equal to $p(c)=\int_0^r \frac{1}{c}
f_2(\frac{t}{c}) (1-\frac{t}{r}) dt$, where $f_2(x)=\frac{2}{\sqrt{2 \pi}} e^{-x^2/2}$.
Therefore
\begin{eqnarray*}
p(c) & = &  \frac{2}{\sqrt{2 \pi}} \int_0^r \frac{1}{c} e^{-(\frac{t}{c})^2/2}  dt  -  \frac{2}{\sqrt{2 \pi}} \int_0^r \frac{1}{c} e^{-(\frac{t}{c})^2/2} \frac{t}{r} dt\\
&      = & S_1(c) - S_2(c)
\end{eqnarray*}

Note that $S_1(c) \le 1$. Moreover
\begin{eqnarray*}
&     S_2(c)=& \frac{2}{\sqrt{2 \pi}} \cdot \frac{c}{r} \int_0^r  e^{-(\frac{t}{c})^2/2} \frac{t}{c^2} dt\\
&     S_2(c) = & \frac{2}{\sqrt{2 \pi}} \cdot  \frac{c}{r} \int_0^{r^2/(2c^2)} e^{-y} dy\\
&     S_2(c) = & \frac{2}{\sqrt{2 \pi}}  \frac{c}{r} [1-e^{-r^2/(2c^2)}]
\end{eqnarray*}

We have $p(1)= S_1(1)-S_2(1) \ge 1-e^{r^2/2} -\frac{2}{\sqrt{2\pi}r} \ge 1-A/r$, for some constant $A>0$.
This implies that the probability that $p$ collides with $q$ is at least $(1-A/r)^k \approx e^{-A}$. 
Thus the algorithm is correct with constant probability.


If $c^2 \le r^2/2$, then we have 
\[  p(c) \le 1- \frac{2}{\sqrt{2 \pi}}\frac{c}{r} (1-1/e) \]
or equivalently $p(c) \le 1-B c/r$, for proper constants $B>0$.



Now consider the expected number of points colliding with $q$.
Let $C$ be a multiset containing all values of $c=\|p'-q\|/\|p-q\|$ over $p' \in P$.
We have
\begin{eqnarray*}
E[|P \cap g^{-1}(q)|] & = & \sum_{c \in C} p(c)^k\\
& = &  \sum_{c \in C, c \le r/\sqrt{2}} p(c)^k +  \sum_{c \in C, c > r/\sqrt{2}} p(c)^k\\
& \le &  \sum_{c \in C, c \le r/\sqrt{2}} (1- B c/r)^k + (1- \frac{B}{\sqrt{2}})^r n\\
%& \le &  _{r/\sqrt{2}}^{\infty} p(c)^k [N(q,c)-N(q,c-dc)] dc\\
& \le  & \int_1^{r/\sqrt{2}} (1-B c/r)^k N(q,c+1) dc + O(1)\\
& \le  & \int_1^{r/\sqrt{2}} e^{-Bc} N(q,c+1) dc + O(1)
\end{eqnarray*}

If $N(q,t)=O(c^b)$, then we have 
\[ E[|P \cap g^{-1}(q)|] =  O \left (\int_1^{r/\sqrt{2}} e^{-Bc} (c+1)^b dc\right ) = 2^{O(b)} \]  


%\hfill $\Box$

\section{Asymptotic Analysis for the General Case}


\begin{theorem}
\label{thm:general_p}
For any $p \in (0,2]$ there is a $(r_1,r_2,p_1,p_2)$-sensitive family
${\cal H}$ for $l_p^d$ such that for any $\gamma > 0$,
$$\rho=\frac{\ln 1/p_1}{\ln 1/p_2} \le (1+\gamma)\cdot\max\left(\frac{1}{c^p}, \frac{1}{c}\right).$$
\end{theorem}


We prove that for the general case ($p \in (0,2]$) the ratio $\rho(c)$
gets arbitrarily close to $\max\left(\frac{1}{c^p},
\frac{1}{c}\right)$.
%Thus, for the case $p=\{1,2\}$, our algorithm matches the best previously
%known $\rho=\frac{1}{c}$.
For the case $p<1$, our algorithm is the
first algorithm to solve this problem, and so there is no existing ratio
against which we can compare our result.  However, we show that for this
case $\rho$ is arbitrarily close to $\frac{1}{c^p}$.
%Note, no algorithm was previously known for the case
%$p <1$, against which we can compare the ratio $\rho(c)$, which for this
%case is arbitrarily close to $\frac{1}{c^p}$.
The proof follows from the
following two Lemmas, which together imply Theorem~\ref{thm:general_p}.

Let $l=\frac{1-p_2}{1-p_1}$, $x=1-p_1$.  Then
$\rho=\frac{\log(1-x)}{\log(1-lx)}\leq\frac{1-p_1}{1-p_2}$ by the
following lemma.

\begin{lemma}
For $x\in[0,1)$ and $l\geq 1$ such that $1 - lx > 0$,
$$\frac{\log(1-x)}{\log(1-lx)}\leq\frac{1}{l}.$$
\end{lemma}

\noindent{\bf Proof:} Noting $\log (1-lx) < 0$, the claim is equivalent to
$l\log(1-x)\geq\log(1-lx)$.
This in turn is equivalent to
%Raising each side of the claim to the power of $2$, we
%see an equivalent statement is
$$g(x)\equiv(1-lx)-(1-x)^l\leq0.$$
This is trivially true for $x=0$.  Furthermore, taking the derivative, we see
$g'(x)=-l+l(1-x)^{l-1}$, which is non-positive for $x\in[0,1)$ and $l\geq1$.
Therefore, $g$ is non-increasing in the region in which we are interested,
and so $g(x)\leq0$ for all values in this region.


Now our goal is to upper bound $\frac{1-p_1}{1-p_2}$.

\begin{lemma}
\label{lemma:general_p}
For any $\gamma >0$, there is $r=r(c,p,\gamma)$ such that
$$\frac{1-p_1}{1-p_2} \le (1+\gamma)\cdot\max\left(\frac{1}{c^p}, \frac{1}{c}\right).$$
\end{lemma}

\noindent{\bf Proof:} Using the values of $p_1, p_2$ calculated in
Sub-section~\ref{subsec:family}, followed by a change of variables, we get

\begin{eqnarray*}
\frac{1-p_1}{1-p_2}
&=&
\frac{1-\int_0^r(1-\frac{t'}{r})f(t')dt'}
     {1-\int_0^r(1-\frac{t'}{r})\frac{1}{c}f(\frac{t'}{c})dt'}
\cr&=&
\frac{1-\int_0^r(1-\frac{t}{r})f(t)dt}
     {1-\int_0^{r/c}(1-\frac{tc}{r})f(t)dt}
\cr&=&
\frac{(1-\int_0^rf(t)dt)+\frac{1}{r}\int_0^rtf(t)dt}
     {(1-\int_0^{r/c}f(t)dt)+\frac{c}{r}\int_0^{r/c}tf(t)dt}.
\end{eqnarray*}
Setting
$$F(x)=1-\int_0^xf(t)dt$$
and
$$G(x)=\frac{1}{x}\int_0^xtf(t)dt$$
we see
\begin{eqnarray*}
\frac{1-p_1}{1-p_2} 
&=& \frac{F(r)+G(r)}{F(r/c)+G(r/c)}\cr
&\leq& \max\left(\frac{F(r)}{F(r/c)},\frac{G(r)}{G(r/c)}\right).
\end{eqnarray*}

First, we consider $p \in (0,2)-\{1\}$ and discuss the special
cases $p=1$ and $p=2$ towards the end.
We bound $F(r)/F(r/c)$.  Notice $F(x)=\Pr_a[a>x]$ for $a$
drawn according to the absolute value of
a $p$-stable distribution with density function
$f(\cdot)$.  To estimate $F(x)$, we can use the Pareto estimation
(\cite{Nol}) for the cumulative distribution function, which holds for $0
< p < 2$,

\begin{tabular}{c}
\hfill~\\
$\forall\delta>0\ \exists x_0\mathrm{\ s.t.\ }\forall x\geq x_0,$\hfill~\\
~~$C_px^{-p}(1-\delta)\leq F(x)\leq C_px^{-p}(1+\delta)$\hfill~\\
\hfill~\\
\end{tabular}

\noindent
where $C_p=\frac{2}{\pi}\Gamma(p)\sin(\pi p/2)$. Note that the extra factor 2 is due to the fact that
the distribution function is for the absolute value of the $p$-stable
distribution. Fix $\delta =
\min(\gamma/4,1/2)$. For this value of $\delta$ let $r_0$ be the $x_0$ in the
equation above.

%Find the $r_0$ corresponding to $\delta\leq\min(\gamma/4,1/2)$ for the
%$\gamma$ in the theorem statement and choose $r\geq r_0$ (notice
%we have the power to set $r$ as we please).  Then

If we set $r>r_0$ we get
\begin{eqnarray*}
\frac{F(r)}{F(r/c)} &\leq&
\frac{C_pr^{-p}(1+\delta)}{C_p(r/c)^{-p}(1-\delta)} \cr&\leq&
\frac{r^{-p}(1+\delta)(1+2\delta)}{(r/c)^{-p}} \cr&\leq&
\left(\frac{1}{c}\right)^p\left(1+4\delta\right) \cr&\leq&
\left(\frac{1}{c}\right)^p\left(1+\gamma\right).
\end{eqnarray*}

Now we bound $G(r)/G(r/c)$.  We break the proof down into two
cases based on the value of $p$.  

\ \\{\bf Case 1:} $p>1$.  For these
$p$-stable distributions, $\int_0^{\infty}tf(t)dt$ converges to,
say, $k_p$ (since the random variables drawn from those distributions have finite expectations).  As $tf(t)$ is non-negative on $[0,\infty)$,
$\int_0^xtf(t)dt$ is a monotonically increasing function of $x$
which converges to $k_p$. Thus, for every $\delta' > 0$ there is some $r_0'$
such that
$$(1-\delta')k_p\leq\int_0^{r_0'}tf(t)dt.$$
Set $\delta'=\min(\gamma/2,1/2)$ and choose $r'>cr_0'$.
Then
\begin{eqnarray}
\frac{G(r')}{G(r'/c)} &=& \frac{\frac{1}{r'}\int_0^{r'}tf(t)dt}
     {\frac{c}{r'}\int_0^{r_0'}tf(t)dt+\frac{c}{r'}\int_{r_0'}^{r'/c}tf(t)dt}
\cr&\leq& \frac{\frac{1}{r'}\int_0^\infty tf(t)dt}
     {\frac{c}{r'}\int_0^{r_0'}tf(t)dt}
\cr&\leq& \frac{\frac{1}{r'}k_p}
     {\frac{c}{r'}(1-\delta')k_p}
\cr&\leq& \frac{1}{c}(1+2\delta') \cr&\leq&
\frac{1}{c}(1+\gamma).
\end{eqnarray}

{\noindent \bf Case 2:} $p<1$. For this case we will choose our parameters
so that we can use the Pareto estimation for the density function.  Choose
$x_0$ large enough so that the Pareto estimation is accurate to within a
factor of $(1\pm\delta)$ for $x>x_0$.  Then for $x>x_0$,

\begin{tabular}{c}
\hfill~\\
$G(x) =
\frac{1}{x}\int_0^{x_0}tf(t)dt+\frac{1}{x}\int_{x_0}^xtf(t)dt$\hfill~\\
\hfill~\\
~~~~~~~$<\frac{1}{x}\int_0^{x_0}tf(t)dt+\frac{1+\delta}{x}\int_{x_0}^xpC_pt^{-p}dt$\hfill~\\
\hfill~\\
~~~~~~~$=\frac{1}{x}\int_0^{x_0}tf(t)dt\ +$\hfill~\\
~~~~~~~~~~$(\frac{pC_p}{x(1-p)}x^{-p+1}-\frac{pC_p}{x(1-p)}{x_0}^{-p+1})(1+\delta)$\hfill~\\
\hfill~\\
~~~~~~~$=\frac{1}{x}(\int_0^{x_0}tf(t)dt-\frac{pC_p(1+\delta)}{(1-p)}{x_0}^{-p+1})\
+$\hfill~\\
~~~~~~~~~~$\frac{1}{x^p}(\frac{pC_p}{(1-p)}(1+\delta)).$\hfill~\\
\hfill~\\
\end{tabular}

%% \begin{eqnarray*}
%% G(x) &=&
%% \frac{1}{x}\int_0^{x_0}tf(t)dt+\frac{1}{x}\int_{x_0}^xtf(t)dt
%% \cr&<&\frac{1}{x}\int_0^{x_0}tf(t)dt+\frac{1+\delta}{x}\int_{x_0}^xpC_pt^{-p}dt
%% \cr&=&
%% \frac{1}{x}\int_0^{x_0}tf(t)dt+\left(\frac{pC_p}{x(1-p)}x^{-p+1}-\frac{pC_p}{x(1-p)}{x_0}^{-p+1}\right)\left(1+\delta\right)
%% \cr&=&
%% \frac{1}{x}\left(\int_0^{x_0}tf(t)dt-\frac{pC_p(1+\delta)}{(1-p)}{x_0}^{-p+1}\right)+\frac{1}{x^p}\left(\frac{pC_p}{(1-p)}(1+\delta)\right).
%% \end{eqnarray*}

\noindent
Since $x_0$ is a constant that depends on $\delta$, the first term
decreases as $1/x$ while the second term decreases as $1/x^p$ where $p<1$.
%the first term of $G(x)$ is inversely linear in $x$
%while the second term is inversely sublinear in $x$.
Thus for every $\delta'$ there is some $x_1$ such that for all $x>x_1$,
the first term is at most $\delta'$ times the second term.  We
choose $\delta'=\delta$. Then for $x>\max(x_1,x_0)$,
$$G(x)<(1+\delta)^2\left(\frac{pC_p}{(1-p)x^p}\right).$$
\iffalse
To lower bound $G(x)$, we again begin with the Pareto estimation.
Using the same parameters as in the upper bound, we see for
$x>x_0$,
\begin{eqnarray*}
G(x) &=&
\frac{1}{x}\int_0^{x_0}tf(t)dt+\frac{1}{x}\int_{x_0}^xtf(t)dt
\cr&>&\frac{1-\delta}{x}\int_{x_0}^xpC_pt^{-p}dt \cr&=&
\left(\frac{pC_p}{x(1-p)}x^{-p+1}-\frac{pC_p}{x(1-p)}{x_0}^{-p+1}\right)(1-\delta)
\cr&=&
\frac{1}{x^p}\left(\frac{pC_p(1-\delta)}{(1-p)}\right)-\frac{1}{x}\left(\frac{pC_p(1-\delta)}{(1-p)}{x_0}^{-p+1}\right)
\end{eqnarray*}
As $p<1$, the first term of $G(x)$ is inversely sublinear in $x$
while the second term is inversely linear in $x$.  Thus for every
$\delta'$ there is some $r_0^2$ such that for all $r>r_0^2$, the
second term is at most $\delta'$ times the first term.  We choose
$\delta'=\delta$.  Then for $r>\max(r_0^2,x_0)$,
\fi
In the same way we obtain
$$G(x)>(1-\delta)^2\left(\frac{pC_p}{(1-p)x^p}\right).$$
Using these two bounds, we see for $r>c\max(x_1,x_0)$,
\begin{eqnarray*}
\frac{G(r)}{G(r/c)} &<&
\frac{(1+\delta)^2\left(\frac{pC_p}{(1-p)r^p}\right)}
     {(1-\delta)^2\left(\frac{pC_p}{(1-p)(r/c)^p}\right)}
\cr&\leq& \frac{1}{c^p}(1+9\delta) \cr&\leq&
\frac{1}{c^p}(1+\gamma)
\end{eqnarray*}
for $\delta<\min(\gamma/9,1/2)$.

We now consider the special cases of $p \in \{1,2\}$.
For the case of $p=1$, we have the Cauchy distribution and we can
compute directly $G(r)={\ln (r^2+1)\over \pi r}$ and $F(r)=1-{2\over
\pi}\tan^{-1} (r)$.
In fact for the ratio $F(r)\over F(r/c)$, the previous analysis for
general $p$ works here.  As for the ratio
${G(r)\over G(r/c)}$, we can prove the upper bound of $1\over c$ using
L'Hopital rule, as follows:

\begin{eqnarray*}
\lim_{r\rightarrow \infty}{G(r)\over G({r\over c})} 
&=& 
\lim_{r\rightarrow \infty}{\ln(r^2+1)\over c\ln((r/c)^2+1)} 
\cr&=& 
\lim_{r\rightarrow \infty}{{2r\over (r^2+1)}\over c({2r\over c^2(r^2/c^2
+1)})}
\cr&=& 
\lim_{r\rightarrow \infty}{c^2(r^2/c^2+1)\over c(r^2+1)} 
\cr&=& 
\lim_{r\rightarrow \infty}{c^2(2r/c^2)\over 2cr} 
\cr&=& {1\over c}.
\end{eqnarray*}

Also for the case $p=2$, i.e. the normal distribution, the computation is
straightforward. We use the fact that for this case $F(r)\simeq f(r)/r$
and $G(r) = {2\over \sqrt{2\pi}}{1-e^{{-{r^2}}\over 2} \over r}$, where
$f(r)$ is the normal density function. For large values of $r$, $G(r)$
clearly dominates $F(r)$, because $F(r)$ decreases exponentially
($e^{-r^2/2}$) while $G(r)$ decreases as $1/r$. Thus, we need to approximate
${G(r)\over G(r/c)}$ as $r$ tends to infinity, which is clearly $1\over c$.

$$\lim_{r\rightarrow \infty}{G(r)\over G({r\over c})} =
\lim_{r\rightarrow \infty} {1-e^{{-{r^2}}\over 2}\over c(1-e^{-{{r^2\over 2c^2}}})} =
{1\over c}$$

Notice that similar to the
previous parts, we can find the appropriate $r(c,p,\gamma)$ such that
$\frac{1-p_1}{1-p_2}$ is at most
most $(1+\gamma){1\over c}$.




\end{document}

\subsection{Asymptotic Analysis for the General Case}
We prove that for the general case ($p \in (0,2]$) the ratio $\rho(c)$
gets arbitrarily close to $\max\left(\frac{1}{c^p},
\frac{1}{c}\right)$. Note, no algorithm was previously known for the case
$p <1$, against which we can compare the ratio $\rho(c)$, which for this
case is arbitrarily close to $\frac{1}{c^p}$. The proof follows from the
following two Lemmas, which together imply Theorem~\ref{thm:general_p}.

Let $l=\frac{1-p_2}{1-p_1}$, $x=1-p_1$.  Then
$\rho=\frac{\log(1-x)}{\log(1-lx)}\leq\frac{1-p_1}{1-p_2}$ by the
following lemma.

\begin{lemma}
For $x\in[0,1)$ and $l\geq 1$ such that $1 - lx > 0$,
$$\frac{\log(1-x)}{\log(1-lx)}\leq\frac{1}{l}.$$
\end{lemma}

\noindent{\bf Proof:} Noting $\log (1-lx) < 0$, the claim is equivalent to
$l\log(1-x)\geq\log(1-lx)$.
This in turn is equivalent to
%Raising each side of the claim to the power of $2$, we
%see an equivalent statement is
$$g(x)\equiv(1-lx)-(1-x)^l\leq0.$$
This is trivially true for $x=0$.  Furthermore, taking derivative, we see
$g'(x)=-l+l(1-x)^{l-1}$, is non-positive for $x\in[0,1)$ and $l\geq1$.
Therefore, $g$ is non-increasing in the region in which we are interested,
and so $g(x)\leq0$ for all values in this region.


Now our goal is to upper bound $\frac{1-p_1}{1-p_2}$.

\begin{lemma}
\label{lemma:general_p}
For an arbitrarily small $\gamma$, there is $r=r(c,p,\gamma)$ such that
$$\frac{1-p_1}{1-p_2} \le (1+\gamma)\cdot\max\left(\frac{1}{c^p}, \frac{1}{c}\right).$$
\end{lemma}

\noindent{\bf Proof:} From the bounds on $p_1$ and $p_2$ calculated above,
along with a change of variables, we see

\begin{eqnarray}
\frac{1-p_1}{1-p_2}
&=&
\frac{1-\int_0^r(1-\frac{t'}{r})f(t')dt'}
     {1-\int_0^r(1-\frac{t'}{r})\frac{1}{c}f(\frac{t'}{c})dt'}
\cr&=&
\frac{1-\int_0^r(1-\frac{t}{r})f(t)dt}
     {1-\int_0^{r/c}(1-\frac{tc}{r})f(t)dt}
\cr&=&
\frac{(1-\int_0^rf(t)dt)+\frac{1}{r}\int_0^rtf(t)dt}
     {(1-\int_0^{r/c}f(t)dt)+\frac{c}{r}\int_0^{r/c}tf(t)dt}.
\end{eqnarray}
Setting
$$F(x)=1-\int_0^xf(t)dt$$
and
$$G(x)=\frac{1}{x}\int_0^xtf(t)dt$$
we see
$$\frac{1-p_1}{1-p_2}\leq\frac{F(r)+G(r)}{F(r/c)+G(r/c)}\leq
\max\left(\frac{F(r)}{F(r/c)},\frac{G(r)}{G(r/c)}\right).$$

First, consider the special cases of $p \in \{1,2\}$.
For the case of $p=1$, we have the Cauchy distribution and we can
compute directly $G(r)={\ln (r+1)\over r}$ and $F(r)=1-{2\over
\pi}\tan^{-1} (r)$ implying $\frac{1-p_1}{1-p_2}$ is at most $1/c$
as $r$ tends to infinity. Also for normal distribution, the
computation is straightforward.  We see $F(r)=f(r)/r$ and $G(r) =
{1-e^{-{r^2}} \over r}$ and so $\frac{1-p_1}{1-p_2}$ is again at
most $1/c$ as $r$ tends to infinity.
{\bf PIOTR: Elaborate ? Also, do we need $r$ going to infinity ?}

It remains to consider $p \in (0,2)-\{1\}$.
We first bound $F(r)/F(r/c)$.  Notice $F(x)=\Pr_a[a>x]$ for $a$
drawn according to a $p$-stable distribution with density function
$f(\cdot)$.  To estimate $F(x)$, we can use the Pareto estimation
(\cite{Nol}) which claims
$$\forall\delta>0\ \exists x_0\mathrm{\ s.t.\ }\forall x\geq x_0,\ C_px^{-p}(1-\delta)\leq F(x)\leq C_px^{-p}(1+\delta)$$
where $C_p=\frac{1}{\pi}\Gamma(p)\frac{\sin(\pi p)}{2}$. Fix $\delta =
\min(\gamma/4,1/2)$. For this value of $\delta$ let $r_0$ be the $x_0$ in
the equation above.
%Find the $r_0$ corresponding to $\delta\leq\min(\gamma/4,1/2)$ for the
%$\gamma$ in the theorem statement and choose $r\geq r_0$ (notice
%we have the power to set $r$ as we please).  Then
If we set $r>r_0$ we get
\begin{eqnarray}
\frac{F(r)}{F(r/c)} &\leq&
\frac{C_pr^{-p}(1+\delta)}{C_p(r/c)^{-p}(1-\delta)} \cr&\leq&
\frac{r^{-p}(1+\delta)(1+2\delta)}{(r/c)^{-p}} \cr&\leq&
\left(\frac{1}{c}\right)^p\left(1+4\delta\right) \cr&\leq&
\left(\frac{1}{c}\right)^p\left(1+\gamma\right).
\end{eqnarray}

Now we bound $G(r)/G(r/c)$.  We break the proof down into two
cases based on the value of $p$.  First suppose $p>1$.  For these
$p$-stable distributions, $\int_0^{\infty}tf(t)dt$ converges to,
say, $k_p$ (since the random variables drawn from those distributions have finite expectations).  As $tf(t)$ is non-negative on $[0,\infty)$,
$\int_0^xtf(t)dt$ is a monotonically increasing function of $x$
which converges to $k_p$. Thus, for every $\delta' > 0$ there is some $r_0'$
such that
$$(1-\delta')k_p\leq\int_0^{r_0'}tf(t)dt.$$
Set $\delta'=\min(\gamma/2,1/2)$ and choose $r'>cr_0'$.
Then
\begin{eqnarray}
\frac{G(r')}{G(r'/c)} &=& \frac{\frac{1}{r'}\int_0^{r'}tf(t)dt}
     {\frac{c}{r'}\int_0^{r_0'}tf(t)dt+\frac{c}{r'}\int_{r_0'}^{r'/c}tf(t)dt}
\cr&\leq& \frac{\frac{1}{r'}\int_0^\infty tf(t)dt}
     {\frac{c}{r'}\int_0^{r_0'}tf(t)dt}
\cr&\leq& \frac{\frac{1}{r'}k_p}
     {\frac{c}{r'}(1-\delta')k_p}
\cr&\leq& \frac{1}{c}(1+2\delta') \cr&\leq&
\frac{1}{c}(1+\gamma).
\end{eqnarray}

If $p<1$, we will choose our parameters so that we can use the
Pareto estimation for the density function.  Choose $x_0$ large
enough so that the Pareto estimation is accurate to within a
factor of $(1\pm\delta)$ for $x>x_0$.  Then for $x>x_0$,
\begin{eqnarray}
G(x) &=&
\frac{1}{x}\int_0^{x_0}tf(t)dt+\frac{1}{x}\int_{x_0}^xtf(t)dt
\cr&<&\frac{1}{x}\int_0^{x_0}tf(t)dt+\frac{1+\delta}{x}\int_{x_0}^xpC_pt^{-p}dt
\cr&=&
\frac{1}{x}\int_0^{x_0}tf(t)dt+\left(\frac{pC_p}{x(1-p)}x^{-p+1}-\frac{pC_p}{x(1-p)}{x_0}^{-p+1}\right)\left(1+\delta\right)
\cr&=&
\frac{1}{x}\left(\int_0^{x_0}tf(t)dt-\frac{pC_p(1+\delta)}{(1-p)}{x_0}^{-p+1}\right)+\frac{1}{x^p}\left(\frac{pC_p}{(1-p)}(1+\delta)\right).
\end{eqnarray}
Since $x_0$ is constant that that depends on $\delta$, the first term
decreases as $1/x$ while the second term decreases as $1/x^p$ where $p<1$.
%the first term of $G(x)$ is inversely linear in $x$
%while the second term is inversely sublinear in $x$.
Thus for every $\delta'$ there is some $x_1$ such that for all $x>x_1$,
the first term is at most $\delta'$ times the second term.  We
choose $\delta'=\delta$. Then for $x>\max(x_1,x_0)$,
$$G(x)<(1+\delta)^2\left(\frac{pC_p}{(1-p)x^p}\right).$$

\iffalse
To lower bound $G(x)$, we again begin with the Pareto estimation.
Using the same parameters as in the upper bound, we see for
$x>x_0$,
\begin{eqnarray}
G(x) &=&
\frac{1}{x}\int_0^{x_0}tf(t)dt+\frac{1}{x}\int_{x_0}^xtf(t)dt
\cr&>&\frac{1-\delta}{x}\int_{x_0}^xpC_pt^{-p}dt \cr&=&
\left(\frac{pC_p}{x(1-p)}x^{-p+1}-\frac{pC_p}{x(1-p)}{x_0}^{-p+1}\right)(1-\delta)
\cr&=&
\frac{1}{x^p}\left(\frac{pC_p(1-\delta)}{(1-p)}\right)-\frac{1}{x}\left(\frac{pC_p(1-\delta)}{(1-p)}{x_0}^{-p+1}\right)
\end{eqnarray}
As $p<1$, the first term of $G(x)$ is inversely sublinear in $x$
while the second term is inversely linear in $x$.  Thus for every
$\delta'$ there is some $r_0^2$ such that for all $r>r_0^2$, the
second term is at most $\delta'$ times the first term.  We choose
$\delta'=\delta$.  Then for $r>\max(r_0^2,x_0)$,
\fi
In the same way we obtain
$$G(x)>(1-\delta)^2\left(\frac{pC_p}{(1-p)x^p}\right).$$
Using these two bounds, we see for $r>c\max(x_1,x_0)$,
\begin{eqnarray}
\frac{G(r)}{G(r/c)} &<&
\frac{(1+\delta)^2\left(\frac{pC_p}{(1-p)r^p}\right)}
     {(1-\delta)^2\left(\frac{pC_p}{(1-p)(r/c)^p}\right)}
\cr&\leq& \frac{1}{c^p}(1+9\delta) \cr&\leq&
\frac{1}{c^p}(1+\gamma)
\end{eqnarray}
for $\delta<\min(\gamma/9,1/2)$.


\section{Conclusions and Open Problems}

In this paper we presented a simple and general algorithm for Approximate
Nearest Neighbor, based on Locality-Sensitive Hashing.
The algorithm appears to be a promising tool for high-performance proximity search in
high dimensions.
Experimental evaluation is under way, and will be reported in the final version of this paper.

It remains open to prove that the optimal ratio $\rho(c)$ is strictly less
than $1/c$ for the case $p=2$. It would be interesting to prove
Lemma~\ref{lemma:general_p} for $r$ values that are bounded.

Achlioptas~\cite{Ach} proved that for the purpose of dimensionality
reduction (under $l_2$ norm) taking the projection of vector with a random
vector whose entries are uniform $\{-1,1\}$ is as efficient as taking the
projection with a random vector whose entries are $N(0,1)$. It is worth
exploring if that technique leads to an efficient LSH scheme.

LaTeX Warning: Label `Fig:n-speedup1' multiply defined.


LaTeX Warning: Label `Fig:n-speedup2' multiply defined.


LaTeX Warning: Label `Fig:n-speedup' multiply defined.

LaTeX Warning: Reference `Fig:eps-speedup1' on page 8 undefined on input line 1
091.


LaTeX Warning: Reference `Fig:eps-speedup2' on page 8 undefined on input line 1
091.


LaTeX Warning: Reference `Fig:mem' on page 8 undefined on input line 1113.
%\addtolength{\baselineskip}{1mm}
\begin{abstract}
We present a novel Locality-Sensitive Hashing scheme for the
decision version of the Approximate Nearest Neighbor Problem
%$(R,c)$-Point Location in Equal Balls (NN) problem\footnote{See
%Preliminaries}
in $l_p^d$ for $p \in (0,2]$, based on $p$-stable distributions.
% and analyze its performance in solving the $\epsilon$-approximate NN problem.
In particular, this gives the first known approximate NN algorithm for the case $p <1$.
Moreover, unlike earlier schemes, our LSH scheme
works directly on points in $\Re^d$ without embeddings. Consequently, the
resulting query time bound is free of large polylogarithmic factors and is
simple and easy to implement. For the important case $p=2$ (Euclidean norm), it
strictly improves the known exponent of $n$ in the space requirement and query
time.
\end{abstract}
